\thispagestyle{empty}

\noindent 
\textsf{\textbf{\Large Master Thesis Abstract}}

\medskip


\noindent
Slovak University of Technology in Bratislava\\
Faculty of Electrical Engineering and Information Technology

\bigskip
\noindent
\begin{tabular}{@{}l l}
Branch of Study: & Cybernetics \\
Study Programme: & Robotics and Cybernetics \\
Author: & Bc. Andrii Stoliar \\
Thesis Title: & Implementation Design of Selected Control Algorithms \\
Supervisor: & Ing. Marián Tárník, PhD. \\
Month, Year: & May, 2026 \\
\end{tabular}

\bigskip
\noindent
Keywords: control algorithm implementation, laboratory dynamic system, system identification, PID controller, adaptive control, MATLAB/Simulink, Arduino UNO, C++, programmable logic controller, PLC

\bigskip
\noindent
This master’s thesis presents the design and implementation of a control system for a laboratory device representing a real dynamic controlled system. Three control algorithms are implemented: a PI controller, a PI controller with operating point switching, and an adaptive controller of the MRAC type. The controlled system is a thermal system in which the controlled variable is the temperature of air forced into a tube. Documentation of the system components and the communication interface is provided for the controlled system. The necessary identification of mathematical models of the system, considering both its static and dynamic properties, is also carried out.
The implementation of the control algorithms is realized on three different platforms: in the MATLAB/Simulink software, using the C++ programming language on an Arduino UNO development board, and by means of an industrial PLC system. For each case, a theoretical analysis of the algorithm, its summary in the form of pseudocode, and technical documentation of the resulting implementation are presented. The thesis also includes verification of each implemented control algorithm on a real laboratory device, together with a comparative evaluation of the results obtained for the same algorithm implemented on different platforms.
The results demonstrate the practical possibilities of implementing the selected algorithms on various platforms while also highlighting the influence of measurement, timing, signal processing, and hardware characteristics on the resulting behavior of the control system.





\pagebreak